%% kapa: exam=1 round=22 year=2011 period=1 form=A subject=civil_law subject_ko=민법 seq=12 number=12
%% source: Q-Net 감정평가사 기출문제 (cst00309) artlSeq=5036790 file=extracted/1008/5036790/2011년 제22회 감정평가사 1차시험 기출문제/1교시.hwp
%% conversion: kordoc 4.24.51 -> tools/gen1.py (auto); review=proofread-sonnet; flags=none
\begin{question}{12}
甲은 자신의 토지를 은닉하기 위해 乙과 짜고 매매를 원인으로 乙 앞으로 소유권이전청구권 보전을 위한 가등기를 해 준 뒤, 다시 丙에게 그 토지를 매도하고 소유권이전등기를 해 주었다. 그 후 가등기에 기하여 乙명의의 본등기가 경료되었고, 乙은 그 토지를 丁에게 매도하고 소유권이전등기를 해 주었다. 丙은 乙명의의 가등기와 본등기가 허위임을 알게 되었다. 다음 설명으로 옳지 않은 것은? (다툼이 있으면 판례에 의함)
\begin{choices}
  \item 甲과 乙 사이의 매매계약 및 가등기는 무효이다.
  \item 乙과 丁 사이의 매매계약은 유효하다.
  \item 乙명의의 본등기를 할 때 등기관은 丙명의의 소유권이전등기를 직권으로 말소한다.
  \item 丁이 甲과 乙 사이의 가장행위에 대하여 알았다면, 丙은 丁에 대하여 말소등기를 구할 수 있다.
  \item 만약 甲과 乙이 짜고 허위로 매매계약을 한 것이 아니라면, 乙은 가등기한 때로부터 그 토지의 소유자이다.
\end{choices}
\end{question}
\answer{5}
